\chapter{Datos secuenciales y redes recurrentes}
\section{Apuntes del Tema 10: Redes Neuronales Recurrentes}

\subsection{1. De Datos Estáticos a Datos Secuenciales}

\subsubsection{1.1 Limitaciones de los Modelos Feedforward (MLPs y CNNs)}

\begin{itemize}

\item \textbf{Asunción de Muestras Independientes (\(i.i.d.\)):} Los modelos tradicionales como los Perceptrones Multicapa (MLP) o las Redes Convolucionales (CNN) asumen que las muestras de entrada son independientes e idénticamente distribuidas.

\end{itemize}

\begin{itemize}

\item \textbf{Entradas de Tamaño Fijo:} Un MLP o una CNN requieren un vector o tensor de dimensiones de entrada rígidamente predefinidas.

\end{itemize}

\begin{itemize}

\item \textbf{Incapacidad para Capturar Dependencias Temporales:} En datos con orden temporal o secuencial (como texto, audio, series temporales financieras o trayectorias GPS), la información del elemento actual depende de los elementos anteriores.

\end{itemize}

\subsubsection{1.2 Modelos Autorregresivos y el Estado Oculto}

Para modelar secuencias \(X = (x_1, x_2, \dots, x_T)\), deseamos estimar la probabilidad conjunta mediante la regla de la cadena de probabilidades:

\$\$P(x\_1, x\_2, \textbackslash\{\}dots, x\_T) = \textbackslash\{\}prod\_\{t=1\}\textasciicircum{}\{T\} P(x\_t \textbackslash\{\}mid x\_1, \textbackslash\{\}dots, x\_\{t-1\})\$\$

\begin{itemize}

\item \textbf{Modelos de Marcov (Orden \(k\)):} Asumen que \(x_t\) solo depende de los \(k\) pasos anteriores.

\end{itemize}

\begin{itemize}

\item **Mecanismo de Estado Oculto (\emph{Hidden State}):\textbf{ En lugar de truncar la historia a \(k\) pasos, una Red Neuronal Recurrente (RNN) mantiene un }estado latente \(h_t\)** que actúa como una memoria comprimida de toda la secuencia procesada desde \(t=1\) hasta \(t-1\).

\end{itemize}

\subsection{2. Arquitectura de la Red Neuronal Recurrente Básica (RNN Vanilla)}

\subsubsection{2.1 Ecuaciones del Estado Oculto y Salida}

En cada paso de tiempo \(t\), la RNN recibe la entrada actual \(x_t \in \mathbb{R}^d\) y el estado oculto del paso anterior \(h_{t-1} \in \mathbb{R}^h\).

\begin{enumerate}

\item \textbf{Actualización del Estado Oculto:}

\end{enumerate}

\$\$h\_t = \textbackslash\{\}phi\textbackslash\{\}left(x\_t W\_\{xh\} + h\_\{t-1\} W\_\{hh\} + b\_h\textbackslash\{\}right)\$\$

\begin{enumerate}

\item \textbf{Cálculo de la Salida:}

\end{enumerate}

\$\$\textbackslash\{\}hat\{y\}\_t = \textbackslash\{\}text\{Softmax\}\textbackslash\{\}left(h\_t W\_\{hq\} + b\_q\textbackslash\{\}right) \textbackslash\{\}quad \textbackslash\{\}text\{o\} \textbackslash\{\}quad o\_t = h\_t W\_\{hq\} + b\_q\$\$

donde:

\begin{itemize}

\item \(W_{xh} \in \mathbb{R}^{d \times h}\): Matriz de pesos para la entrada.

\end{itemize}

\begin{itemize}

\item \(W_{hh} \in \mathbb{R}^{h \times h}\): Matriz de pesos para la transición recurrente de estados.

\end{itemize}

\begin{itemize}

\item \(W_{hq} \in \mathbb{R}^{h \times q}\): Matriz de pesos para la capa de salida.

\end{itemize}

\begin{itemize}

\item \(\phi(\cdot)\): Función de activación no lineal (típicamente \(\tanh\) o \(\text{ReLU}\)).

\end{itemize}

\begin{quote}**Compartición de Pesos (\emph{Weight Sharing} Temporal):** Las matrices de pesos \(\{W_{xh}, W_{hh}, W_{hq}\}\) son exactamente las mismas en todos los pasos de tiempo \(t=1, \dots, T\).\end{quote}

\begin{quote}\end{quote}

\begin{quote}\end{quote}

\begin{lstlisting}[language=Python]

Estructura Desplegada en el Tiempo (Unfolded RNN):



  y_1         y_2                 y_T

   ▲           ▲                   ▲

   │           │                   │

┌─────┐     ┌─────┐             ┌─────┐

│ h_1 │───► │ h_2 │───► ... ───►│ h_T │

└─────┘ W_hh└─────┘ W_hh        └─────┘

   ▲           ▲                   ▲

   │ W_xh      │ W_xh              │ W_xh

  x_1         x_2                 x_T



\end{lstlisting}

\subsection{3. Entrenamiento: Retropropagación a través del Tiempo (BPTT)}

\subsubsection{3.1 Backpropagation Through Time (BPTT)}

Para entrenar las RNNs, la red recurrente se "despliega" a lo largo de los \(T\) pasos temporales formando un grafo acíclico dirigido. La pérdida total es la suma de las pérdidas en cada paso de tiempo:

\$\$\textbackslash\{\}mathcal\{L\}(W) = \textbackslash\{\}sum\_\{t=1\}\textasciicircum{}\{T\} \textbackslash\{\}mathcal\{L\}\_t(y\_t, \textbackslash\{\}hat\{y\}\_t)\$\$

Para calcular el gradiente respecto a la matriz de transición recurrente \(W_{hh}\), se aplica la regla de la cadena acumulando las dependencias pasadas:

\$\$\textbackslash\{\}frac\{\textbackslash\{\}partial \textbackslash\{\}mathcal\{L\}\_t\}\{\textbackslash\{\}partial W\_\{hh\}\} = \textbackslash\{\}sum\_\{k=1\}\textasciicircum{}\{t\} \textbackslash\{\}frac\{\textbackslash\{\}partial \textbackslash\{\}mathcal\{L\}\_t\}\{\textbackslash\{\}partial h\_t\} \textbackslash\{\}cdot \textbackslash\{\}left( \textbackslash\{\}prod\_\{j=k+1\}\textasciicircum{}\{t\} \textbackslash\{\}frac\{\textbackslash\{\}partial h\_j\}\{\textbackslash\{\}partial h\_\{j-1\}\} \textbackslash\{\}right) \textbackslash\{\}cdot \textbackslash\{\}frac\{\textbackslash\{\}partial h\_k\}\{\textbackslash\{\}partial W\_\{hh\}\}\$\$

\subsubsection{3.2 Problema Crítico: Desvanecimiento y Explosión del Gradiente}

El término \(\prod_{j=k+1}^{t} \frac{\partial h_j}{\partial h_{j-1}}\) implica multiplicar repetidamente la matriz de pesos \(W_{hh}^\top\).

\begin{itemize}

\item **Desvanecimiento del Gradiente (\emph{Vanishing Gradients}):\textbf{ Si los autovalores de \(W_{hh}\) son menores a \(1\), la magnitud del gradiente decrece exponencialmente hacia cero cuando la secuencia es larga (\(T \gg 0\)). La red }olvida las dependencias a largo plazo**.

\end{itemize}

\begin{itemize}

\item **Explosión del Gradiente (\emph{Exploding Gradients}):\textbf{ Si los autovalores son mayores a \(1\), el gradiente crece exponencialmente destruyendo el entrenamiento. Se soluciona mediante }Gradient Clipping**:

\end{itemize}

\$\$g \textbackslash\{\}leftarrow \textbackslash\{\}min\textbackslash\{\}left(1, \textbackslash\{\}frac\{\textbackslash\{\}theta\}\{\textbackslash\{\}Vert\{\}g\textbackslash\{\}Vert\{\}\}\textbackslash\{\}right) g\$\$

\subsection{4. Redes Gated: LSTM y GRU}

Para solucionar el desvanecimiento de gradiente y permitir recordar información a largo plazo, se crearon arquitecturas basadas en puertas de control (\emph{gates}).

\subsubsection{4.1 LSTM (\emph{Long Short-Term Memory})}

La LSTM sustituye la celda estándar por una estructura de memoria que mantiene un \textbf{estado de celda \(C_t\)} (\emph{Cell State} o autopista de información) y un \textbf{estado oculto \(h_t\)}.

\begin{lstlisting}[language=Python]

             ┌──────────────────────────────────────────────┐ (Cell State C_t)

             │                 x (f_t)        + (i_t * C~_t)│

C_{t-1} ─────┴─────────────────┴──────────────┴─────────────┴────────► C_t

                   ▲                ▲         ▲          ▲

                   │                │         │          │

h_{t-1} ─┬─────────┼───────────────┼─────────┼──────────┼─────────┬──► h_t

         │         │                │         │          │         │

x_t ─────┴─► [Forget Gate]    [Input Gate] [Candidate] [Output Gate]



\end{lstlisting}

\subsubsection{Puertas y Ecuaciones Matemáticas de la LSTM:}

\begin{enumerate}

\item \textbf{Forget Gate (\(f_t\)):} Decide qué información descartar del estado de celda anterior.

\end{enumerate}

\$\$f\_t = \textbackslash\{\}sigma\textbackslash\{\}left(x\_t W\_\{xf\} + h\_\{t-1\} W\_\{hf\} + b\_f\textbackslash\{\}right)\$\$

\begin{enumerate}

\item \textbf{Input Gate (\(i_t\)) y Candidato (\(\tilde{C}_t\)):} Deciden qué nueva información almacenar.

\end{enumerate}

\$\$i\_t = \textbackslash\{\}sigma\textbackslash\{\}left(x\_t W\_\{xi\} + h\_\{t-1\} W\_\{hi\} + b\_i\textbackslash\{\}right)\$\$

\$\$\textbackslash\{\}tilde\{C\}\_t = \textbackslash\{\}tanh\textbackslash\{\}left(x\_t W\_\{xc\} + h\_\{t-1\} W\_\{hc\} + b\_c\textbackslash\{\}right)\$\$

\begin{enumerate}

\item \textbf{Actualización del Cell State (\(C_t\)):}

\end{enumerate}

\$\$C\_t = f\_t \textbackslash\{\}odot C\_\{t-1\} + i\_t \textbackslash\{\}odot \textbackslash\{\}tilde\{C\}\_t\$\$

\begin{enumerate}

\item \textbf{Output Gate (\(o_t\)) y Estado Oculto (\(h_t\)):}

\end{enumerate}

\$\$o\_t = \textbackslash\{\}sigma\textbackslash\{\}left(x\_t W\_\{xo\} + h\_\{t-1\} W\_\{ho\} + b\_o\textbackslash\{\}right)\$\$

\$\$h\_t = o\_t \textbackslash\{\}odot \textbackslash\{\}tanh(C\_t)\$\$

\subsubsection{4.2 GRU (\emph{Gated Recurrent Unit})}

Variante simplificada de la LSTM que combina el estado de celda y el estado oculto en una sola variable \(h_t\), reduciendo el coste computacional y el número de parámetros.

\begin{enumerate}

\item \textbf{Reset Gate (\(r_t\)):} Determina cuánta información del pasado olvidar:

\end{enumerate}

\$\$r\_t = \textbackslash\{\}sigma\textbackslash\{\}left(x\_t W\_\{xr\} + h\_\{t-1\} W\_\{hr\} + b\_r\textbackslash\{\}right)\$\$

\begin{enumerate}

\item \textbf{Update Gate (\(z_t\)):} Actúa simultáneamente como puerta de olvida e entrada:

\end{enumerate}

\$\$z\_t = \textbackslash\{\}sigma\textbackslash\{\}left(x\_t W\_\{xz\} + h\_\{t-1\} W\_\{hz\} + b\_z\textbackslash\{\}right)\$\$

\begin{enumerate}

\item \textbf{Candidato de Estado Oculto (\(\tilde{h}_t\)):}

\end{enumerate}

\$\$\textbackslash\{\}tilde\{h\}\_t = \textbackslash\{\}tanh\textbackslash\{\}left(x\_t W\_\{xh\} + (r\_t \textbackslash\{\}odot h\_\{t-1\}) W\_\{hh\} + b\_h\textbackslash\{\}right)\$\$

\begin{enumerate}

\item \textbf{Estado Oculto Final (\(h_t\)):}

\end{enumerate}

\$\$h\_t = (1 - z\_t) \textbackslash\{\}odot h\_\{t-1\} + z\_t \textbackslash\{\}odot \textbackslash\{\}tilde\{h\}\_t\$\$

\subsection{5. RNNs Bidireccionales y Apiladas (\emph{Deep RNNs})}

\begin{itemize}

\item \textbf{RNN Bidireccional (BiRNN / BiLSTM):} Procesa la secuencia en ambas direcciones simultáneamente usando dos estados ocultos independientes (\(\overrightarrow{h}_t\) hacia adelante y \(\overleftarrow{h}_t\) hacia atrás). El estado combinado es \(h_t = [\overrightarrow{h}_t \,;\, \overleftarrow{h}_t]\).

\end{itemize}

\begin{itemize}

\item \emph{Uso:} Crucial en tareas de procesamiento de lenguaje natural (NLP) donde el contexto futuro influye en el significado del presente.

\end{itemize}

\begin{itemize}

\item **RNN Profunda / Apilada (\emph{Stacked RNN}):** Apila múltiples capas recurrentes donde la salida del estado oculto de la capa \(l\) sirve como entrada \(x_t\) para la capa \(l+1\).

\end{itemize}

\subsection{6. Práctica y Aplicación en Código (PyTorch)}

A continuación se incluye una implementación modular completa que ilustra la construcción de una red \textbf{LSTM Bidireccional} para la clasificación de secuencias (ej. análisis de sentimiento o series temporales) en PyTorch:

\begin{lstlisting}[language=Python]

import torch

import torch.nn as nn

import torch.optim as optim



# --------------------------------------------------------------------------

# 1. Definición del Modelo BiLSTM / GRU en PyTorch

# --------------------------------------------------------------------------

class SequenceClassifier(nn.Module):

    def __init__(self, vocab_size=1000, embed_dim=64, hidden_dim=128, num_classes=2, num_layers=2, dropout=0.3):

        super(SequenceClassifier, self).__init__()

        

        # Capa de Embedding para convertir índices de tokens en vectores continuos

        self.embedding = nn.Embedding(vocab_size, embed_dim)

        

        # Capa LSTM Bidireccional de múltiples capas

        self.lstm = nn.LSTM(

            input_size=embed_dim,

            hidden_size=hidden_dim,

            num_layers=num_layers,

            batch_first=True,       # Entrada esperada: (Batch, Seq_Len, Embed_Dim)

            bidirectional=True,     # BiLSTM

            dropout=dropout if num_layers > 1 else 0.0

        )

        

        # Capa lineal final: hidden_dim * 2 por la bidireccionalidad

        self.fc = nn.Linear(hidden_dim * 2, num_classes)

        self.dropout = nn.Dropout(p=dropout)



    def forward(self, x):

        # x shape: (Batch, Seq_Len)

        embedded = self.embedding(x) # (Batch, Seq_Len, Embed_Dim)

        

        # output shape: (Batch, Seq_Len, Hidden_Dim * 2)

        # h_n, c_n shape: (Num_Layers * 2, Batch, Hidden_Dim)

        lstm_out, (h_n, c_n) = self.lstm(embedded)

        

        # Concatenar el último estado oculto forward y backward de la última capa

        hidden_forward = h_n[-2, :, :]

        hidden_backward = h_n[-1, :, :]

        last_hidden = torch.cat((hidden_forward, hidden_backward), dim=1) # (Batch, Hidden_Dim * 2)

        

        out = self.dropout(last_hidden)

        logits = self.fc(out) # (Batch, Num_Classes)

        return logits



# --------------------------------------------------------------------------

# 2. Ejemplo de Bucle de Entrenamiento

# --------------------------------------------------------------------------

def train_rnn_demo():

    device = torch.device("cuda" if torch.cuda.is_available() else "cpu")

    

    # Hiperparámetros

    VOCAB_SIZE = 1000

    BATCH_SIZE = 16

    SEQ_LEN = 32

    NUM_CLASSES = 2

    

    # Modelo, pérdida y optimizador

    model = SequenceClassifier(

        vocab_size=VOCAB_SIZE, 

        embed_dim=64, 

        hidden_dim=128, 

        num_classes=NUM_CLASSES, 

        num_layers=2

    ).to(device)

    

    criterion = nn.CrossEntropyLoss()

    optimizer = optim.Adam(model.parameters(), lr=0.001)

    

    # Datos sintéticos de secuencia: (Batch, Seq_Len)

    dummy_input = torch.randint(0, VOCAB_SIZE, (BATCH_SIZE, SEQ_LEN)).to(device)

    dummy_labels = torch.randint(0, NUM_CLASSES, (BATCH_SIZE,)).to(device)

    

    model.train()

    optimizer.zero_grad()

    

    # Forward Pass

    logits = model(dummy_input)

    loss = criterion(logits, dummy_labels)

    

    # Backward Pass

    loss.backward()

    

    # Gradient Clipping esencial para evitar la explosión de gradiente en RNNs

    nn.utils.clip_grad_norm_(model.parameters(), max_norm=5.0)

    

    optimizer.step()

    

    print(f"Demostración RNN completada en dispositivo: {device}")

    print(f"Dimensión de salida (Logits): {logits.shape}")

    print(f"Pérdida inicial (Loss): {loss.item():.4f}")



if __name__ == "__main__":

    train_rnn_demo()



\end{lstlisting}
\section{Conexión con las prácticas}

Los notebooks pasan del modelado de secuencias con MLP a una RNN implementada desde cero para predicción del siguiente carácter y exponen la necesidad de controlar el gradiente.

\subsection{Fuente práctica: \texttt{10\_01\_SequentialData.ipynb}}

Ruta: code/10\_01\_SequentialData/10\_01\_SequentialData.ipynb

\textbf{Apartados detectados:} Character-level Language Modeling with MLPs.

Salidas gráficas almacenadas: 0.

\textbf{Salida textual conservada:} \begin{quote}'malena es un cuy' | --- Starting training: OneHotMLP --- | Epoch [1/10] - Avg Loss: 2.6257 | Epoch [2/10] - Avg Loss: 2.2111\end{quote}

\textbf{Fragmento representativo:}

\begin{lstlisting}[language=Python]

import os
import math
from collections import Counter

import numpy as np
import pandas as pd
import matplotlib.pyplot as plt
import seaborn as sns

import torch
import torch.nn as nn
import torch.nn.functional as F

import torch.optim as optim
from torch.utils.data import Dataset
from torch.utils.data import DataLoader

\end{lstlisting}

\subsection{Fuente práctica: \texttt{10\_02\_RNNfromScratch.ipynb}}

Ruta: code/10\_02\_RNN/10\_02\_RNNfromScratch.ipynb

\textbf{Apartados detectados:} Implementing a Recurrent Neural Network from scratch for next-character prediction.

Salidas gráficas almacenadas: 0.

\textbf{Salida textual conservada:} \begin{quote}Device: cuda | Number of characters: 24937 | Walt Whitman has somewhere a fine and just distinction between “loving | by allowance” and “loving with personal love.” This distinction applies\end{quote}

\textbf{Fragmento representativo:}

\begin{lstlisting}[language=Python]

import torch
import torch.nn as nn
from torch.utils.data import Dataset, DataLoader
import math
import random

\end{lstlisting}
\section{Autoevaluación razonada}

\subsection*{Pregunta} ¿Qué información conserva el estado oculto?

\textbf{Solución.} Una representación del contexto procesado hasta el paso temporal actual.

\subsection*{Pregunta} ¿Qué es gradient clipping?

\textbf{Solución.} Una limitación de la norma del gradiente para evitar actualizaciones excesivas, especialmente útil frente a explosión del gradiente.
\section{Fuentes y trazabilidad}

\textbf{Diapositivas:}

\begin{itemize}

\item \path{slides/10\_01\_SequentialData.pdf}

\item \path{slides/10\_02\_RNN.pdf}

\end{itemize}

\textbf{Nota personal:} apuntes/Apuntes del Tema 10 - Redes Neuronales Recurrentes.md